%% =====================================================================
%%  B4-T-22-01 -- what B4 contributes to "The Liking Handle"
%%  -------------------------------------------------------------------
%%  LAYER A: APPEND-ONLY. Written once, then left alone. Material found
%%  only here sits in the `unique` slot and is protected by rule R10.
%% =====================================================================
%% @kind: source
%% @id: B4-T-22-01
%% @book: B4
%% @topic: T-22-01
%% @locator: ch. 5, pp. 157--177
%% @status: distilled
%% @unique: yes
%% @integrated: yes
%% @updated: 2026-09-19

\dossier{B4}{T-22-01}{\bookshort{B4}}{ch. 5, pp. 157--177}

\begin{distilled}
  \item Compliance rises with how much the requester is liked, and the effect is independent of the strength of the case.
  \item Liking is produced by resemblance, by praise, by familiarity developed through cooperation, by surface appeal, and by association with things already valued.
  \item Each component can be manufactured without the quality it signals.
  \item Praise retains its effect on a recipient who suspects it is strategic.
\end{distilled}

\begin{unique}
  \item The decomposition of liking into five producible components. B1 treats favouritism as a tactic and B3 as courtiership; neither supplies the list of cues that generate the feeling.
  \item The finding that praise works on a recipient who knows it is praise, which is what makes the handle durable rather than fragile.
\end{unique}
